\documentclass[11pt]{article}
\usepackage[a4paper,margin=1.04in]{geometry}
\usepackage[T1]{fontenc}
\usepackage{lmodern,amsmath,amssymb,amsthm,mathtools,mathrsfs,microtype}
\usepackage{enumitem,booktabs,array}
\usepackage[hidelinks]{hyperref}
\hypersetup{pdftitle={Factorial Moments in Two Variables: Exact Leaf Transport and Trace Orthogonality},pdfauthor={Christopher D. Long}}
\setlength{\emergencystretch}{2em}
\numberwithin{equation}{section}
\newtheorem{theorem}{Theorem}[section]
\newtheorem{lemma}[theorem]{Lemma}
\newtheorem{proposition}[theorem]{Proposition}
\newtheorem{corollary}[theorem]{Corollary}
\theoremstyle{definition}\newtheorem{definition}[theorem]{Definition}
\theoremstyle{remark}\newtheorem{remark}[theorem]{Remark}
\newcommand{\C}{\mathbb C}
\newcommand{\R}{\mathbb R}
\newcommand{\Q}{\mathbb Q}
\newcommand{\Qb}{\overline{\mathbb Q}}
\newcommand{\Lfac}{\mathcal L}
\newcommand{\CC}{\mathcal C}
\newcommand{\HH}{\widehat G}
\newcommand{\dd}{\,d}
\newcommand{\Crit}{\operatorname{CritVal}}
\newcommand{\Span}{\operatorname{span}}
\newcommand{\Res}{\operatorname{Res}}
\newcommand{\ord}{\operatorname{ord}}
\newcommand{\UC}{U_d^{\C}}
\newcommand{\GL}{\operatorname{GL}}
\newcommand{\one}{\mathbf 1}
\title{Factorial Moments in Two Variables:\\
Exact Leaf Transport and Trace Orthogonality}
\author{Christopher D. Long\\\small\texttt{galizur@gmail.com}}
\date{7 October 2026}
\begin{document}
\maketitle
\begin{abstract}
Let $\Lfac(x^ay^b)=a!b!$. We prove factorial-moment nonvanishing for
nonzero polynomials in two variables whose highest homogeneous part is
not a scalar multiple of $(x+y)^D$. The proof uses an exact radial
Mittag--Leffler transform, polynomial monodromy, and
Lindemann--Weierstrass. The transform produces vanishing leaf germs
containing a pole density and a Volterra correction. The corrections
continue with their full integration histories retained and have no
negative Puiseux coefficients at a nonzero terminal critical value.
The resulting weighted coefficient annihilates the top rational
monodromy component, while the unweighted inverse-derivative vector
belongs to that component by a zero-trace identity on every proper
composition block. Conjugate pairing excludes all nonzero critical
values, and endpoint leaf germs dispose of the remaining pure powers.
Combining the nonflat theorem with the explicitly stated
single-integral nonvanishing theorem of EIT~I gives the two-variable
Factorial Conjecture. As a formal consequence, the zero-integral
polynomials for the product-exponential measure
$e^{-x-y}\dd x\dd y$ form a Mathieu--Zhao subspace; indeed the
corresponding eventual-power radical is trivial. Thus the
two-dimensional product-exponential case of van den Essen's Integral
Conjecture follows. The nonflat argument is independent of EIT~I.
\end{abstract}

\section{The statement and the inputs}\label{sec:introduction}
Define the complex-linear functional $\Lfac:\C[x,y]\to\C$ by
$\Lfac(x^ay^b)=a!b!$ for $a,b\ge0$. The Factorial Conjecture was introduced
by van den Essen, Wright and Zhao~\cite{FC}. Liu and Sun proved the
homogeneous two-variable case~\cite{LiuSun}. The present argument treats
arbitrary lower homogeneous terms.

For $0\ne f\in\C[x,y]$ of total degree $D\ge1$, write
\begin{equation}\label{eq:radial-polynomial}
 P(t,r)=f(rt,r(1-t))=\sum_{j=0}^{D}A_j(t)r^{D-j},
 \qquad \deg A_j\le D-j,\qquad A=A_0.
\end{equation}
Restriction of degree-$D$ homogeneous forms to $x+y=1$ is injective:
off $x+y=0$, homogeneity recovers the form from this restriction.
Thus $A(t)=f_D(t,1-t)$ is nonzero. Its degree $d$ can be smaller than $D$.
We call the cases $d>0$ and $d=0$ \emph{nonflat} and \emph{flat}, respectively.
In particular,
\[
 A\text{ is constant}\quad\Longleftrightarrow\quad
 f_D(x,y)=c(x+y)^D\quad\text{for some }c\ne0.
\]

\begin{theorem}[Nonflat nonvanishing]\label{thm:nonflat}
If $0\ne f\in\C[x,y]$ has nonconstant $A$, then $\Lfac(f^n)\ne0$ for some
$n\ge1$, and in fact for infinitely many $n$.
\end{theorem}

\begin{theorem}[EIT~I: the external nonvanishing input]\label{thm:eit}
Let $q\in\Qb[t]$ be nonconstant, let $a,b\in\Qb$ be distinct, and let
$\xi\in\Qb^\times$. Then
\[
 \int_a^b e^{\xi q(t)}\dd t\ne0.
\]
\end{theorem}
This is the nonvanishing part of the single-integral corollary of
EIT~I~\cite[Corollary~1.4, source label \texttt{cor:single}]{EIT};
that result also asserts
transcendence. The integral can be taken along any finite path, since its
integrand is entire. No extension to arbitrary amplitudes is being used.

\begin{corollary}[FC(2), using Theorem~\ref{thm:eit}]\label{cor:fc2}
For every $f\in\C[x,y]$,
\[
 \bigl(\Lfac(f^n)=0\text{ for every }n\ge1\bigr)\quad\Longrightarrow\quad f=0.
\]
Equivalently, every nonzero $f$ has infinitely many nonzero factorial moments.
\end{corollary}

The following consequence belongs to the original Image/Mathieu framework
in which the Factorial Conjecture was introduced. Recall that a linear
subspace $M$ of a commutative algebra $R$ is a Mathieu--Zhao subspace if,
whenever $f^m\in M$ for every $m\ge1$, one has $gf^m\in M$ for all
sufficiently large $m$ for every $g\in R$. Van den Essen's Integral
Conjecture predicts this property for the kernel of integration against a
positive measure with finite polynomial moments~\cite{AmazingImage}.
The connection is therefore not new; the new input here is
Corollary~\ref{cor:fc2}.

\begin{corollary}[Two-dimensional product-exponential Integral Conjecture]
\label{cor:integral-mathieu}
Let
\[
 B=(0,\infty)^2,\qquad
 \dd\sigma=e^{-x-y}\dd x\dd y,
\]
and let
\[
 M=\left\{h\in\C[x,y]:
       \int_B h\,\dd\sigma=0\right\}.
\]
Then $M=\ker\Lfac$ is a Mathieu--Zhao subspace of $\C[x,y]$. In fact,
\[
 \left\{f\in\C[x,y]:
   \Lfac(f^m)=0\text{ for all sufficiently large }m\right\}=\{0\}.
\]
Consequently, the Integral Conjecture holds for the product-exponential
measure $e^{-x-y}\dd x\dd y$ in dimension two.
\end{corollary}

\begin{proof}
For monomials,
\[
 \int_B x^a y^b e^{-x-y}\dd x\dd y=a!b!,
\]
so $M=\ker\Lfac$. Suppose first that
$\Lfac(f^m)=0$ for all $m\ge N$. Then for every $k\ge1$,
\[
 \Lfac\bigl((f^N)^k\bigr)=\Lfac(f^{Nk})=0.
\]
Corollary~\ref{cor:fc2} applied to $f^N$ gives $f^N=0$, hence $f=0$.
This proves the displayed eventual-power assertion. In particular, if
$f^m\in M$ for every $m\ge1$, then $f=0$, and therefore
$gf^m=0\in M$ for every $g\in\C[x,y]$ and every $m\ge1$.
Thus $M$ is a Mathieu--Zhao subspace.
\end{proof}

\paragraph{Direct dependencies and scope.}
The nonflat proof uses Lindemann--Weierstrass and the top-component
irreducibility theorem for rational permutation modules of groups
containing a full cycle~\cite[Theorem~3.1 and Section~4.1]{PM}.
These are published inputs. We state precisely the part
of that representation theorem needed below, and prove the endpoint
averaging, planar separation, and resulting path-module containment in
our notation. We do not reprove the underlying Schur-ring classification.
For the unrestricted flat conclusion, the additional external input is
Theorem~\ref{thm:eit}. Its source is a separate research manuscript;
its proof is not reproduced here. The $E$-function lifting and
Fourier--Laplace theory used there are transitive dependencies of that
input. Apart from standard complex analysis, algebraic inverse branches,
and the Nullstellensatz and L\"uroth theorems, the remaining main-text
arguments are included. None of the appendices is a logical input to
either main conclusion.
\paragraph{Manuscript status.}
This is an unrefereed research manuscript, not a Lean formalization.
The EIT dependence of Corollary~\ref{cor:fc2} is explicit; the nonflat
argument is independent of EIT. The exact checks accompanying the source
verify particular identities, not the universal theorem.

\section{Normalization and the exact kernel}\label{sec:kernel}
\begin{lemma}[Algebraic specialization]\label{lem:specialization}
A complex counterexample with specified total degree $D$ and specified
positive degree $d$ of $A$ would give a counterexample over $\Qb$ with the
same degrees. The same assertion holds with the normalized flat top $A=1$.
\end{lemma}
\begin{proof}
Use the finitely many coefficients of a polynomial of degree at most $D$
as variables. All equations $\Lfac(f^n)=0$ have rational coefficients.
Impose the linear conditions specifying $\deg A=d$, and adjoin the inverse
of its degree-$d$ coefficient. In the flat case impose $A=1$, after
scaling the proposed counterexample by its nonzero constant leading value.
Evaluation at the assumed complex point shows that the ideal over
$\Qb$ is proper. Noetherianity and the Nullstellensatz give an algebraic
point annihilating it. Every moment polynomial belongs to the ideal, so
this point satisfies the entire family of moment equations. The inverse
equation ensures that the degree-$d$ coefficient of $A$ is nonzero;
hence $\deg A=d$ and $f_D\ne0$, preserving total degree $D$.
In the flat case $A=1$ likewise ensures $f_D\ne0$.
Neither the monodromy group nor the decomposition lattice has to survive
specialization, since the contradiction is uniform in both.
\end{proof}

Set
\begin{align}
 M_n&=\frac{\Lfac(f^n)}{\Gamma(Dn+2)}, &
 \HH(w)&=\frac1w\sum_{n\ge0}M_nw^{-n},\label{eq:normalized}\\
 E_{D,2}(z)&=\sum_{n\ge0}\frac{z^n}{\Gamma(Dn+2)}, &
 \CC(w,t)&=\frac1w\int_0^\infty e^{-r}r
 E_{D,2}\!\left(\frac{P(t,r)}w\right)\dd r.\label{eq:kernel}
\end{align}
Termwise integration of a polynomial gives
\[
 \Lfac(f^n)=\int_0^\infty\!\int_0^\infty e^{-x-y}f(x,y)^n\dd x\dd y
 =\int_0^1\!\int_0^\infty e^{-r}rP(t,r)^n\dd r\dd t.
\]
The last equality uses the Jacobian $r$ of $x=rt$, $y=r(1-t)$.
For sufficiently large $|w|$ the resulting generating integral is normally
convergent, and
\begin{equation}\label{eq:full-transform}
 \HH(w)=\int_0^1\CC(w,t)\dd t.
\end{equation}
Indeed, for $x=|z|^{1/D}$ the defining series gives the global bound
\begin{equation}\label{eq:ML-growth}
 |E_{D,2}(z)|\le\sum_{n\ge0}\frac{x^{Dn}}{(Dn+1)!}
 \le\sum_{m\ge0}\frac{x^m}{m!}=e^x.
\end{equation}
Since $|P(t,r)|^{1/D}\le C(1+r)$ on a compact $t$-set, this estimate
also justifies summation under the integral for sufficiently large $|w|$.
All positive moments vanish exactly when $\HH(w)=1/w$ near infinity.

Put $V=\{t:A(t)\ne0\}$. We need a coefficient on a finite cover, not a
single-valued logarithmic coefficient on the entire base.
\begin{lemma}[Exact local kernel and a continued coefficient]\label{lem:kernel}
Near $w=A(t)\ne0$, a continued germ of the kernel has the form
\begin{equation}\label{eq:local-kernel}
 \CC(w,t)=\frac{E(t)}{w-A(t)}+
 B_0(w,t)\log(w-A(t))+H(w,t),\qquad
 E(t)=\exp\!\left(\frac{A_1(t)}{DA(t)}\right),
\end{equation}
where $B_0,H$ are locally jointly holomorphic. The residue is unchanged
by continuation. For fixed $t\in V$, the only possible finite exceptional
values for continuation in $w$ are $0,A(t)$, locally uniformly in $t$.

The distinguished diagonal germ of $B_0$ extends to a holomorphic function
$\mathscr B$ on
\begin{equation}\label{eq:cover}
 \mathscr X=\{(w,t,\rho):w\ne0,\ t\in V,\ \rho^D=A(t)/w\}.
\end{equation}
It is given by
\begin{equation}\label{eq:global-B}
 \mathscr B(w,t,\rho)=\frac{\rho^{D-1}}{DA(t)}
 \sum_{m\ge1}\frac{(-1)^mF_m(t)}{(m-1)!}(1-\rho)^{m-1},
\end{equation}
where $F_m$ are defined in \eqref{eq:Laurent}. The series is entire in
$1-\rho$, locally uniformly on $V$. The initial diagonal germ is the
restriction with $\rho$ near $1$.
\end{lemma}
\begin{proof}
Let $\omega=e^{2\pi i/D}$. Expansion of the exponentials gives the exact
root filter
\begin{equation}\label{eq:filter}
 E_{D,2}(X^D)=\frac1{DX}\sum_{k=0}^{D-1}\omega^{-k}e^{\omega^kX}
 -C_{D,1}(X),\quad
 C_{D,1}=\begin{cases}X^{-1},&D=1,\\0,&D\ge2.\end{cases}
\end{equation}
Fix a compact set $K\subset V$. Split the radial integral at a
sufficiently large $R$, chosen uniformly for $t\in K$, and apply the
filter only to the exterior part. Leave the compact part in its entire
Mittag--Leffler form. Initially impose $|w|>|A(t)|$, so
$|\rho|=|A(t)/w|^{1/D}<1$ and all $D$ exponential summands converge.
Each summand is then continued separately from this common domain.

On $V$, the exterior root with leading term $r$ and its inverse are
\[
 \Phi_t(r)=(P(t,r)/A(t))^{1/D}=r+p_1(t)+O(r^{-1}),\quad
 \chi(t,Y)=Y-p_1(t)+O(Y^{-1}),\quad p_1=\frac{A_1}{DA}.
\]
The analytic inverse theorem at infinity gives convergent Laurent
series on common exterior disks for $t\in K$; one first works on a
neighborhood of $K$ and then uses compactness. The chosen leading terms
specify the roots and inverses uniquely, so their coefficients are
single-valued and holomorphic on $V$. Define
\begin{equation}\label{eq:Laurent}
 F(t,Y)=e^{Y-\chi(t,Y)}\frac{\chi(t,Y)\chi_Y(t,Y)}Y
       =\sum_{m\ge0}F_m(t)Y^{-m}.
\end{equation}
Here $F_0=E$. By a uniform Cauchy estimate at infinity, there are
$C_K,R_0>0$ such that $|F_m(t)|\le C_KR_0^m$ for every $t\in K$ and
$m\ge0$. Enlarge $R$ so that the image paths lie in this common
exterior domain.
With $\rho=(A(t)/w)^{1/D}$ near $1$, the summand singular at this diagonal is
\begin{equation}\label{eq:singular-summand}
 \frac{\rho^{D-1}}{DA(t)}
 \int_{\Phi_t([R,\infty))}e^{-(1-\rho)Y}F(t,Y)\dd Y.
\end{equation}
In $\Re(1-\rho)>0$ the exterior image ray can be moved to a common positive
ray. The far connector tends to zero exponentially, and the finite
connector is entire in $1-\rho$, locally uniformly in $t$.

Choose a starting point $R_1>R_0$ on that ray and put
\[
 J_m(\eta)=\int_{R_1}^\infty e^{-\eta Y}Y^{-m}\dd Y.
\]
Direct integration and integration by parts give
\begin{align*}
 J_0&=e^{-\eta R_1}/\eta,\\
 J_1&=-\log\eta-\gamma-\log R_1-
       \sum_{j\ge1}\frac{(-\eta R_1)^j}{j\,j!},\\
 J_{m+1}&=\frac{e^{-\eta R_1}R_1^{-m}-\eta J_m}{m}\qquad(m\ge1).
\end{align*}
The logarithmic coefficient of $J_m$ is
$(-1)^m\eta^{m-1}/(m-1)!$. For $m\ge1$ put
\[
 H_m(\eta)=J_m(\eta)
 -\frac{(-1)^m\eta^{m-1}}{(m-1)!}\log\eta.
\]
The displayed formulas show that $H_m$ is entire and obeys the same
recurrence as $J_m$. On $|\eta|\le L$, set
$B_m=R_1^{m-1}\sup|H_m|$. Then
\[
 B_{m+1}\le\frac{e^{LR_1}}m+\frac{LR_1}m B_m.
\]
For $m\ge2LR_1$ the last coefficient is at most $1/2$; enlarging a
constant to cover the finitely many earlier terms proves
$\sup_{|\eta|\le L}|H_m|\le C_{L,R_1}R_1^{1-m}$.
Together with $R_1>R_0$ and the bound on $F_m$, this proves normal
convergence of the entire remainders. The logarithmic coefficients
converge normally by their factorial denominators. Thus the tail is
\[
 \frac{E(t)}\eta+
 \left(\sum_{m\ge1}\frac{(-1)^mF_m(t)}{(m-1)!}\eta^{m-1}\right)\log\eta
 +\text{a function holomorphic near }\eta=0.
\]
This supplies exact continuation on the universal cover of $\eta\ne0$.
Since
\[
 1-\rho=\frac{w-A(t)}{DA(t)}+O((w-A(t))^2),
\]
\eqref{eq:singular-summand} has residue $E(t)$ and logarithmic coefficient
\eqref{eq:global-B}. Changing $\log(1-\rho)$ to $\log(w-A(t))$ changes
only a holomorphic remainder.

The entireness used below also has a contour formulation. For any
$R_2>R_0$, orient $|Y|=R_2$ counterclockwise. Termwise integration gives
\[
 \Lambda(t,\eta):=\sum_{m\ge1}\frac{(-1)^mF_m(t)}{(m-1)!}\eta^{m-1}
 =-\frac1{2\pi i}\int_{|Y|=R_2}e^{-\eta Y}F(t,Y)\dd Y.
\]
Thus $|\Lambda(t,\eta)|\le C_{K,R_2}e^{R_2|\eta|}$, uniformly on $K$,
and $\mathscr B=\rho^{D-1}\Lambda(t,1-\rho)/(DA(t))$ is holomorphic
on the stated cover. This contour formula is for $\Lambda$; the
prefactor must be retained when recovering $\mathscr B$.

The other root summands have parameters $1-\omega^k\rho$; their zeros all
project to $w=A(t)$. At $\rho=1$ they are locally holomorphic for $k\ne0$.
The compact part and the $D=1$ correction are holomorphic at every
nonzero diagonal. This proves the claimed exceptional set and uniformity.
A change of logarithm branch adds a function holomorphic on the cover,
and cannot alter the pole residue.
\end{proof}

\begin{remark}[Canonical local coefficient versus continued germ]\label{rem:B-sheets}
The local logarithmic coefficient of the complete kernel is canonical at
its nonzero diagonal. The continuation of its distinguished germ, however,
is $\mathscr B(w,t,\rho)$ on $\mathscr X$ and need not descend to the base
$(w,t)$-space. If another root summand becomes singular, the identity
$\omega^{-k}\rho^{D-1}=(\omega^k\rho)^{D-1}$ gives the same canonical
local coefficient after reselecting that summand. It does not identify
different continued root sheets. Every integral below retains its root
parameter. Appendix~\ref{app:examples} gives an exact sign-changing example.
\end{remark}

\begin{lemma}[Normality near zeros of the top]\label{lem:normality}
Near a zero of $A$, the original kernel integral is normally convergent
for $w$ on any fixed compact set away from zero. On every bounded $t$-disk
it is holomorphic in $t$ for all sufficiently large $|w|$.
\end{lemma}
\begin{proof}
For $D>1$, locally uniformly,
\[
 |P(t,r)/w|^{1/D}\le |A(t)/w|^{1/D}r+C(1+r)^{1-1/D}.
\]
For $D=1$ the second term is bounded. Near a zero of $A$, or for $|w|$
sufficiently large on a bounded $t$-disk, make the first coefficient
less than $1/2$. By \eqref{eq:ML-growth}, on compact $w$-sets away
from zero the kernel integrand is bounded by
\[
 C_1r\exp\bigl(-r/2+C_2(1+r)^{1-1/D}\bigr),
\]
with the last contribution bounded when $D=1$. This is integrable and
gives joint holomorphy locally in $(w,t)$, including at zeros of $A$.
To compare with a continued germ, first fix the entire finite
continuation path in $w$ away from zero. Its modulus has a positive
minimum, so near a zero of $A$ one can shrink the $t$-neighborhood until
$|A(t)/w|^{1/D}<1/2$ along the whole path. The single-valued integral
agrees with the germ near infinity, hence with its continuation along
this path by uniqueness. No uniform claim over an unbounded $t$-set is
needed.
\end{proof}

\section{Leaf identities and continuation with retained paths}\label{sec:transport}
Choose a regular value $b$ different from $0,A(0),A(1)$. Let $T$ be an
embedded finite star centered at $b$, with one arm to each distinct member
of
\[
 \{0,A(0),A(1)\}\cup\Crit(A).
\]
The arms have disjoint interiors and are straight near their leaves.
The inverse graph $\Gamma=A^{-1}(T)$ is a tree. Here is a useful direct
justification. If $T$ has $s$ arms, then $\Gamma$ has $ds$ edges, $d$
center vertices, and $ds-(d-1)$ leaf vertices: the total finite
ramification of a degree-$d$ polynomial is $d-1$. It is connected because
the local monodromies about the leaf values generate a transitive
covering group; their cycles join precisely the incident lifted stars.
Thus its Euler characteristic is $1$, and it is a tree. This is the usual
polynomial inverse-tree construction~\cite[Section~2.1]{PM}.

Let $\mu$ be its unique simple path from the actual point $0$ to the actual
point $1$. Both $\mu$ and $[0,1]$ lie in a bounded disk. Lemma~\ref{lem:normality}
therefore permits replacing the interval by $\mu$ in
\eqref{eq:full-transform} for large $w$.
Append a simple arc from $b$ to infinity, disjoint from $T$ except at
$b$ and straight sufficiently near infinity. Label the inverse branches
over $b$ so that the geometric monodromy obtained by following this arc,
making one counterclockwise turn on a large circle, and returning is
$g_\infty=(1\ 2\ \cdots\ d)$. We retain this geometric labeling throughout;
an arbitrary generator of the same cyclic subgroup need not preserve
the circular order used below. Orient each target arm from its leaf
$v$ toward $b$. The path gives a vector
$n_v\in\{0,1,-1\}^d$, recording the signed use of the lift in each sheet.
With these orientations,
\begin{equation}\label{eq:tree}
 \HH(w)=\sum_v\int_v^b\sum_{i=1}^d n_{v,i}
             \CC(w,t_i(u))t_i'(u)\dd u.
\end{equation}
The complement $\widehat{\C}\setminus T$ in the Riemann sphere is
simply connected. For $t\in\mu$, the possible kernel singularities
$0,A(t)$ lie in $T$, and the kernel is holomorphic at infinity.
Lemmas~\ref{lem:kernel} and~\ref{lem:normality} therefore determine its
single-valued continuation off $T$, locally uniformly along the compact
path $\mu$. Consequently \eqref{eq:tree} holds there by continuation from
infinity. It retains every lift over each leaf.

Fix $v\ne0$ and write $c_i=\lim_{u\to v}t_i(u)$ on its arm. For $w$
near $v$, define the germ
\begin{equation}\label{eq:leaf-germ}
 \mathscr L_v(w)=\sum_i n_{v,i}
 \left\{E(t_i(w))t_i'(w)+
 \int_{c_i}^{t_i(w)}\mathscr B(w,t,\rho)\dd t\right\}.
\end{equation}
Initially the integrals are local near $c_i$. Their lifts to $\mathscr X_w$
are selected by $\rho=1$ at the upper endpoint. Since $A(c_i)=v\ne0$,
regular and ramified endpoints are both allowed.

\begin{lemma}[Complete local singular extraction]\label{lem:local-extraction}
If $\HH$ is rational, then $\mathscr L_v\equiv0$ for every nonzero leaf $v$.
\end{lemma}
\begin{proof}
Only meromorphy of the chosen continued germ of $\HH$ at $v$ is
used in this local argument. Put $z=w-v$ and $x=u-v$. The germ arriving
from the complement of $T$ uses the diagonal root $\rho\to1$; its local
logarithm may differ from another choice only by a holomorphic term.
Near the leaf, set
\begin{align*}
 a(x)&=\sum_i n_{v,i}E(t_i(v+x))t_i'(v+x),\\
 b_0(z,x)&=\sum_i n_{v,i}
 \mathscr B(v+z,t_i(v+x),\rho)t_i'(v+x),
 \qquad \rho^D=(v+x)/(v+z),\quad\rho\sim1.
\end{align*}
These have convergent Puiseux expansions in $x$ with a common finite
denominator; $b_0$ is also holomorphic in $z$. Every $x$-exponent is
greater than $-1$, since a local inverse of degree $k$ has derivative
$O(x^{1/k-1})$.

All other edges in \eqref{eq:tree} are holomorphic in $w$ near $v$ on the
chosen continuation. This follows from Lemma~\ref{lem:kernel} where
$A(t)\ne0$, and from Lemma~\ref{lem:normality} near zero top values.
The integrated local holomorphic remainder is holomorphic as well.
After straightening the short arm, its relevant part is
\[
 \int_0^\epsilon
 \left(\frac{a(x)}{z-x}+b_0(z,x)\log(z-x)\right)\dd x.
\]
A common nonzero linear rescaling of $x,z$ does not affect the following
vanishing argument. Take $z$ in a disk of radius less than $\epsilon$,
slit along its positive radius. Then $z-x\ne0$ for $0\le x\le\epsilon$,
and $\log(z-x)$ has a continuous choice on this product, holomorphic
in $z$; at $x=\epsilon$ it is holomorphic on the full smaller disk.
Put $q(z,x)=\int_0^x b_0(z,s)\dd s$. This tends to zero at the lower endpoint.
Integration by parts gives a holomorphic boundary term plus
\begin{equation}\label{eq:corrected-cauchy}
 \int_0^\epsilon\frac{a(x)+q(z,x)}{z-x}\dd x.
\end{equation}
In particular the sign of the correction is positive.

There are finitely many fractional classes. Choose their representatives
$-1<\delta\le0$ and write
\[
 a(x)+q(z,x)=\sum_\delta x^\delta h_\delta(z,x),
\]
where each $h_\delta$ is holomorphic in the ordinary variables $z,x$.
Divide
$h_\delta(z,x)=h_\delta(z,z)+(x-z)k_\delta(z,x)$.
The $k_\delta$-term in \eqref{eq:corrected-cauchy} is holomorphic in $z$,
by integrability of $x^\delta$. The only nonholomorphic term from this
class is therefore
\[
 h_\delta(z,z)I_\delta(z),\qquad
 I_\delta(z)=\int_0^\epsilon\frac{x^\delta}{z-x}\dd x.
\]
For $-1<\delta<0$, the singular part of $I_\delta$ is a nonzero multiple
of $z^\delta$ on a fixed slit; scaling and the beta integral give a factor
proportional to $\pi/\sin(\pi\delta)$. For $\delta=0$ it is $\log z$,
up to a holomorphic function. Let $T_0$ denote continuation once around
$z=0$, and choose $N$ divisible by all Puiseux denominators. Applying
$T_0^N-1$ to a meromorphic germ kills the fractional terms and isolates
$2\pi iN h_0(z,z)$, so $h_0(z,z)=0$. The remaining terms
$z^\delta h_\delta(z,z)$ have distinct monodromy eigenvalues
$e^{2\pi i\delta}\ne1$. The corresponding polynomial projectors in
$T_0$ isolate each term from the others and from the meromorphic part.
Since the constants multiplying $z^\delta$ are nonzero, every
$h_\delta(z,z)$ vanishes. Hence
\[
 a(z)+\int_0^z b_0(z,x)\dd x=0.
\]
The change of variables $t=t_i(v+x)$ identifies this with
\eqref{eq:leaf-germ}. This proves vanishing of the whole corrected germ,
not only of its leading coefficient.
\end{proof}

\begin{remark}[The triangular coefficient equations]\label{rem:triangular}
If $a(x)=\sum a_\alpha x^\alpha$ and
$b_0(z,x)=\sum b_{p,\beta}z^p x^\beta$, the preceding identity is equivalent to
\begin{equation}\label{eq:triangular}
 a_\alpha+\sum_{p+\beta+1=\alpha}\frac{b_{p,\beta}}{\beta+1}=0
 \qquad(\alpha>-1).
\end{equation}
It is not an identity for $a$ alone. For example,
$-\sqrt{x}/(z-x)+(2\sqrt{x})^{-1}\log(z-x)$ is
$\partial_x(\sqrt{x}\log(z-x))$, and its integral is holomorphic near
$z=0$, although both densities are nonzero.
\end{remark}

Write $G_A$ for the inverse-branch monodromy group. Fix a terminal path
from $b$ toward a nonzero critical value $\lambda$, with a final straight
segment, avoiding zero and all critical values before its endpoint.
For any negative Puiseux exponent $\alpha$, define the linear functional
\begin{equation}\label{eq:coefficient-functional}
 \Phi_{\lambda,\alpha}(e_i)=
 [(w-\lambda)^\alpha]
 \frac{\exp(A_1(t_i(w))/(Dw))}{A'(t_i(w))}
\end{equation}
using continuation along that path and a common Puiseux parameter.

\begin{theorem}[Transport of the full leaf identity]\label{thm:transport}
If $\HH$ is rational, $\Phi_{\lambda,\alpha}$ annihilates
$\Q[G_A]n_v$ for every nonzero leaf $v$ and every $\alpha<0$.
\end{theorem}
\begin{proof}
Let $U=\C\setminus(\{0\}\cup\Crit(A))$. Continue the zero germ
$\mathscr L_v$ along its arm to $b$, and then along any path in $U$.
The initial local integration paths lie in $V$, and every moving
inverse branch stays there because $A(t_i(w))=w\ne0$.
For an individual integral in \eqref{eq:leaf-germ}, retain its path
from the fixed $c_i$ and append
the path of its moving upper endpoint. Lift the accumulated path to
$\mathscr X_w$ with $\rho=1$ at that upper endpoint. The fixed lower
root is allowed to change after a loop; it remains part of the data.

Here is a precise local description. On a small $w$-disk $W$ away from
zero choose a holomorphic root $s(w)^D=w$, and put
\[
 \Sigma=\{(t,\sigma):t\in V,\ \sigma^D=A(t)\},
 \qquad \rho=\sigma/s(w).
\]
Then $\mathscr X|_W$ is identified with $W\times\Sigma$.
The curve $\Sigma$ is locally parametrized by $t$, since $\sigma\ne0$,
even at a critical point of $A$. The differential
$\mathscr B(w,t,\sigma/s(w))\dd t$ is holomorphic on this product.
A fixed compact lifted path integrates to a holomorphic function of $w$;
a short appended piece has a jointly holomorphic local primitive.
This constructs the continuation along each compact portion of the
chosen path. At the moving endpoint $\sigma=s(w)$, so $\rho=1$ is
preserved. Different roots at a fixed lower endpoint are never reset.

Now approach $\lambda\ne0$, and suppose an endpoint tends to $c$ with
$A(c)=\lambda$. Split the accumulated history into a fixed compact lifted
path ending near $c$, and a short terminal piece. The first part is
holomorphic on a disk about $\lambda$ by the product description. A
point with $A(t)=w$ on that history causes no singularity of
$\mathscr B$, because its coefficient series is entire in $1-\rho$.
No extension across $A(t)=0$ is asserted or used.
Near $(\lambda,c)$ the root with $\rho\sim1$ is holomorphic. A primitive
$J(w,t)$ of the corresponding differential is jointly holomorphic, and
the terminal piece is
\[
 J(w,t_i(w))-J(w,t_*).
\]
If $A$ has local degree $k$ at $c$, then
$t_i(w)=c+O((w-\lambda)^{1/k})$. Thus this primitive has only
nonnegative Puiseux exponents. The same holds for every retained history.
Every negative coefficient of the continued identity comes solely from
$E(t_i(w))t_i'(w)$.

For each $g\in G_A$, choose a loop at $b$ realizing $g$, and then follow
the fixed terminal path. The continued $\mathscr L_v$ is still zero, so
its negative coefficient is $\Phi_{\lambda,\alpha}(g n_v)=0$.
Linear extension proves the theorem. Removing the extra regular point
zero from the base, when zero is not critical, does not remove any
permutation from its monodromy image. The complete histories can depend
on more than $g$; no action of $G_A$ on those histories is asserted.
Only the negative Puiseux coefficients factor through the permutation
action: the history terms contribute none, while the pole contribution
depends only on $g$.
\end{proof}

\section{The path module and the simultaneous trace identity}\label{sec:module}
We use the bilinear pairing $\langle x,y\rangle=\sum_i x_i y_i$.
Vector spaces are over $\Q$ until explicitly complexified.
For the actual signed vectors in \eqref{eq:tree}, put
\[
 M=\sum_v\Q[G_A]n_v,\qquad M_{\ne0}=\sum_{v\ne0}\Q[G_A]n_v.
\]
\begin{lemma}[The zero leaf is algebraically redundant]\label{lem:zero-leaf}
One has $\sum_v n_v=0$ coordinatewise, and $M_{\ne0}=M$.
\end{lemma}
\begin{proof}
Each visited center vertex of $\mu$ has one entering and one leaving
edge, contributing opposite signs to its sheet coordinate. Neither actual
endpoint lies over $b$. Thus $n_0=-\sum_{v\ne0}n_v$; apply all monodromy
elements. This makes no analytic assertion at $w=0$.
\end{proof}

Use the geometric labeling of Section~\ref{sec:transport}, with
$g_\infty=(1\ 2\ \cdots\ d)$, and identify sheet $i$ with $i-1$ in
$\mathbb Z/d\mathbb Z$. For $k\mid d$, let $V_k\subset\Q^d$ consist of $k$-periodic vectors, and set
\begin{equation}\label{eq:top-component}
 D(G_A)=\{k\mid d:V_k\text{ is }G_A\text{-invariant}\},\qquad
 U_d=\bigcap_{\substack{k\in D(G_A)\\k<d}}V_k^\perp,
 \qquad \UC=U_d\otimes_\Q\C.
\end{equation}
The indicators of residue classes modulo $k$ span $V_k$. Because the
cyclic subgroup is regular, every block system is of this form: the
block containing the label $0$ in $\mathbb Z/d\mathbb Z$ is a subgroup,
and its translates are all blocks. Invariance of its indicator span is
equivalent to invariance of the partition. In polynomial monodromy these
blocks correspond to decompositions $A=P\circ R$. The intermediate-field
correspondence and L\"uroth's theorem give a rational generator; the
unique point over infinity permits choosing that generator $R$ with
its only pole at infinity. Then $R$, and also $P$, are polynomials.
For $V_k$ one has $\deg P=k$ and $\deg R=d/k$.

\begin{theorem}[Rational cyclic-module input, Pakovich--Muzychuk]
\label{thm:rational-input}
For a permutation group $G\le S_d$ containing a full $d$-cycle, the
space $U_d$ defined as in \eqref{eq:top-component} is an irreducible
rational $G$-module.
\end{theorem}
This is the top-component case of~\cite[Theorem~3.1 and Section~4.1]{PM}.
All numbering in references to~\cite{PM} refers to arXiv:0710.4085v2.
Theorem~3.1 there uses the maximal proper elements below $d$ in $D(G)$.
Intersecting over these or over all proper elements gives the same
space: every proper $k$ divides some maximal proper $f$, and
$V_k\subseteq V_f$. Section~4.1 explicitly identifies the latter
intersection with the irreducible top rational component.
Its proof uses rational Schur rings over a cyclic group. That
representation-theoretic theorem, not the whole polynomial moment problem,
is retained as an external input. The next lemmas establish its application
to the actual path, following~\cite[Section~2.2, Propositions~2.4, 2.6
and~4.1]{PM}.

\begin{lemma}[A rational invariant subspace dichotomy]\label{lem:dichotomy}
If $N\subset\Q^d$ is $G_A$-invariant, then either $U_d\subset N$ or
$N\perp U_d$.
\end{lemma}
\begin{proof}
The orthogonal projectors $P_N$ and $P_U$ have rational entries and commute
with $G_A$. They therefore commute with the full cycle. Its centralizer
consists of circulant matrices and is commutative, so $P_NP_U=P_UP_N$.
The restriction of $P_N$ to $U_d$ is a $G_A$-equivariant idempotent.
By Theorem~\ref{thm:rational-input} it is either zero or the identity.
These are the two assertions. No absolute irreducibility is used.
\end{proof}

For an actual endpoint $p\in\{0,1\}$, let $I_p$ be the set of sheet labels
whose inverse branches on its arm tend to $p$, and let $e_p=|I_p|$ be
its local degree. Write
\[
 m_p=\frac1{e_p}\sum_{i\in I_p}e_i\in\Q^d.
\]
\begin{lemma}[Endpoint averaging]\label{lem:averaging}
If $A(0)=A(1)$, then $m_0-m_1\in M$. If $A(0)\ne A(1)$, then
$m_0,m_1\in M$.
\end{lemma}
\begin{proof}
Average the vector $n_v$ over the finite cyclic group generated by the
local monodromy at its leaf. On the labels tending to one point above
$v$, the averaging replaces every coordinate by the mean over that
local cycle. Equivalently, the average of each basis vector is the
uniform vector on its local-monodromy orbit, even when the monodromy
has several cycles of different lengths. At an interior visit of $\mu$,
one edge enters and another leaves the point; their opposite signs
therefore cancel after averaging. At an actual endpoint
the sum is $+1$ for $0$, and $-1$ for $1$, under our orientations.
Thus if the endpoint values agree, the average is $m_0-m_1$.
Otherwise the two separate averages are $m_0$ and $-m_1$.
All averages belong to $M$.
\end{proof}

\begin{lemma}[Planar separation of endpoint labels]\label{lem:separation}
Let $\zeta=e^{2\pi i/d}$, and identify a label $i$ with $\zeta^{i-1}$.
If $A(0)=A(1)$, the two endpoint label sets lie in disjoint open arcs
of the circle. If their values differ, the sets have this property or
share one point and, after that point is removed, lie on opposite open
arcs cut out by it and one further circle point. In particular, their
convex hulls are disjoint in the first case and meet at most in that
single circle point in the second.
\end{lemma}
\begin{proof}
Use the appended arc already chosen in Section~\ref{sec:transport}.
Its lifts are disjoint rays from the $d$ center vertices to infinity.
Because the arc is straight near infinity, its lifted ends have
asymptotic directions in the geometric cyclic order and each crosses
a sufficiently large circle just once. Truncate there. The resulting
boundary order is the order of $g_\infty$, up to reversal, which does
not affect separation.
For an endpoint $p$, take the union $K_p$ of its incident arm lifts and
the attached rays. It is an embedded star joining $p$ to exactly the
boundary labels $I_p$; intermediate center vertices have degree two.

If the endpoints have the same value, $K_0$ and $K_1$ are disjoint.
Between two consecutive rays of $K_0$, the two rays and the intervening
boundary arc enclose a sector by the Jordan curve theorem. The connected
star $K_1$, disjoint from $K_0$, lies in a single such sector. Its boundary
labels consequently occur in one gap between consecutive labels of
$I_0$. Choose one chord endpoint between the first label of $I_1$ in
this gap and its preceding $I_0$ label, and the other between the last
$I_1$ label and its following $I_0$ label. This chord strictly separates
the two finite sets. If one set is a singleton, a sufficiently small
arc about that label gives the same conclusion.

For different endpoint values, two stars can share a center only if
one lifted star contains both endpoints. There cannot be two such
centers: the two paths from $0$ through these centers to $1$ would
form a cycle in $A^{-1}(T)$. If there is no common center, the preceding
argument applies. Otherwise $K_0\cap K_1$ is the one common ray from that
center to the boundary. After this common ray, including its center,
is removed, the rest of $K_1$ is connected, disjoint from $K_0$, and
has that center in its closure. It lies in one sector of $K_0$
adjacent to the common ray. If both label sets have other labels,
traverse this sector from the common label through the labels of $I_1$.
Join the common label to a point in the gap between the last such
$I_1$ label and the next $I_0$ label. The remaining labels are strictly
on opposite sides of this chord; each convex hull touches it only at
the common label. If either set has no other labels, its hull is that
single point and the separation assertion is immediate. Thus the two
hulls meet at most at the common label in all cases.
\end{proof}

\begin{proposition}[The actual path contains the top component]
\label{prop:path-module}
For $d\ge2$, $U_d\subset M$.
\end{proposition}
\begin{proof}
Put $z=(1,\zeta,\ldots,\zeta^{d-1})$. The sum of its entries on every
residue class modulo a proper divisor $k$ is zero. Thus $z\in\UC$.
The number $\langle m_p,z\rangle$ is the arithmetic mean of the roots
of unity labeled by $I_p$.
If the endpoint values agree, Lemma~\ref{lem:separation} gives disjoint
convex hulls, so these two means differ. Lemma~\ref{lem:averaging}
therefore supplies a vector in $M$ pairing nontrivially with $z$.
If the values differ, both means cannot be zero: their convex hulls
are disjoint or meet only on the unit circle. At least one of the two
vectors in Lemma~\ref{lem:averaging} again pairs nontrivially with $z$.
Orthogonality of the rational spaces would imply orthogonality of
their complexifications, so the complex witness $z$ suffices.
Hence $M$ is not orthogonal to $U_d$, and Lemma~\ref{lem:dichotomy}
gives $U_d\subset M$.
\end{proof}
This is the needed path-module statement of~\cite[Proposition~4.1]{PM},
with its endpoint and sign conventions established here. It allows
ramified endpoints and equal or zero endpoint values. A regular extra
leaf at zero contributes no active lift unless it is an endpoint value.

\begin{corollary}\label{cor:annihilation}
If $\HH$ is rational, every functional \eqref{eq:coefficient-functional}
with $\lambda\ne0$ and $\alpha<0$ annihilates $\UC$.
\end{corollary}
\begin{proof}
Theorem~\ref{thm:transport} gives annihilation of $M_{\ne0}$.
Lemma~\ref{lem:zero-leaf} and Proposition~\ref{prop:path-module} give
$U_d\subset M=M_{\ne0}$. Extend scalars.
\end{proof}

\begin{lemma}[Simultaneous trace orthogonality]\label{lem:trace}
For every regular $w$, in the continued sheet ordering,
\begin{equation}\label{eq:inverse-derivative}
 u(w)=\left(\frac1{A'(t_1(w))},\ldots,\frac1{A'(t_d(w))}\right)\in\UC.
\end{equation}
Every Puiseux coefficient vector of $u$ also belongs to $\UC$.
\end{lemma}
\begin{proof}
Take any proper $k\in D(G_A)$ and the corresponding decomposition
$A=P\circ R$, with $r=\deg R=d/k\ge2$. A block over $w$ is a fiber
$R^{-1}(\xi)$ with $P(\xi)=w$. Regularity gives $P'(\xi)\ne0$ and
$r$ distinct points in that fiber. The block trace is
\begin{equation}\label{eq:block-trace}
 \sum_{i\text{ in block}}\frac1{A'(t_i(w))}
 =\frac1{P'(\xi)}\sum_{R(c)=\xi}\frac1{R'(c)}=0.
\end{equation}
Indeed, in the partial fractions
\[
 \frac1{R(t)-\xi}=\sum_{R(c)=\xi}\frac1{R'(c)(t-c)},
\]
the coefficient of $t^{-1}$ is zero since the left side is $O(t^{-r})$.
Block indicators span $V_k$, so $u\perp V_k$. Intersect over all proper
$k$. The relations are constant rational equations, and consequently
hold coefficientwise in every continued Puiseux expansion.
This uses the rational block-indicator equations, not any reality
assumption on the coefficients of $A$.
\end{proof}
The intersection in this argument treats all proper decompositions at
once. It imposes no constancy condition on a weighted vector along any
one chosen block system.

\section{The arithmetic obstruction and the nonflat conclusion}\label{sec:arithmetic}
The next proposition applies to the top two layers of any algebraically
specialized polynomial; it does not impose genericity on their coefficients.

\begin{proposition}[Nonzero critical values are detected]\label{prop:arithmetic}
Suppose $A,A_1\in\Qb[t]$, $\deg A=d\ge2$, and all negative coefficient
functionals in \eqref{eq:coefficient-functional} at nonzero critical
values annihilate $U_d$. Then $A$ has no nonzero critical value.
\end{proposition}
\begin{proof}
Suppose $\lambda\ne0$ is a critical value. Let $K\ge2$ be the largest
local degree of $A$ among the points above $\lambda$, and put
$\alpha=1/K-1<0$. At a point $c$ of local degree $K$, write
\[
 A(c+s)=\lambda+a_cs^K+O(s^{K+1}),\qquad a_c\ne0.
\]
The $K$ inverse branches there have convergent expansions
\[
 t_i(w)=c+b_c\zeta_i(w-\lambda)^{1/K}
              +O((w-\lambda)^{2/K}),
 \qquad b_c^K=a_c^{-1},\quad \zeta_i^K=1.
\]
Let $\beta$ be the coefficient vector at exponent $\alpha$ in
$u(w)$ from \eqref{eq:inverse-derivative}. On these branches,
$\beta_i=b_c\zeta_i/K$. At a point of smaller local degree $k<K$,
the first exponent $1/k-1$ is strictly larger than $\alpha$, so
$\beta_i=0$. In particular, $\beta\ne0$, its entries are algebraic,
and Lemma~\ref{lem:trace} gives $\beta\in\UC$.

Since $A(t_i(w))=w$ and $\lambda\ne0$, the exponential weight is
holomorphic in the local inverse coordinate. Its constant term is
$\exp(A_1(c_i)/(D\lambda))$, and all other terms have positive
Puiseux exponent. Thus the coefficient at $\alpha$ in the weighted
inverse-derivative vector is exactly
\begin{equation}\label{eq:weighted-leading}
 \phi_i=\begin{cases}
 \beta_i\exp(A_1(c_i)/(D\lambda)),&\beta_i\ne0,\\
 0,&\beta_i=0.
 \end{cases}
\end{equation}
The maximal choice of $K$ is what excludes contributions from smaller
ramification indices or nonconstant terms of the weight.

By hypothesis $\phi$ is bilinearly orthogonal to $\UC$. Since $U_d$
is defined by rational linear equations, its complexification is
preserved by coordinatewise conjugation; $A$ itself need not have
real coefficients. Consequently $\overline\beta\in\UC$, and
\begin{equation}\label{eq:positive-pairing}
 0=\langle\phi,\overline\beta\rangle
  =\sum_{\beta_i\ne0}|\beta_i|^2
       \exp\!\left(\frac{A_1(c_i)}{D\lambda}\right).
\end{equation}
All exponents and coefficients in this expression are algebraic.
After equal exponents are grouped, each coefficient is a strictly
positive real algebraic number. At least one group occurs.
Lindemann--Weierstrass~\cite{Baker} makes the exponentials of distinct
algebraic numbers linearly independent over $\Qb$, contradicting
\eqref{eq:positive-pairing}.
\end{proof}

\begin{remark}
Grouping by critical points instead gives
\[
 \sum_{\substack{A(c)=\lambda\\e_c=K}}
 \frac{(a_c\overline{a_c})^{-1/K}}{K}
 \exp\!\left(\frac{A_1(c)}{D\lambda}\right)=0,
\]
where the positive real root is intended. This records explicitly that
local root choices do not affect the positive coefficients. The
positivity is created by the conjugate pairing, not by the original
oriented interval or by the complex exponentials.
\end{remark}

\begin{lemma}[Only zero critical values]\label{lem:pure-power}
If all critical values of a polynomial $A$ of degree $d\ge2$ are zero,
then $A(t)=a(t-c)^d$ for some $a\ne0$.
\end{lemma}
\begin{proof}
Every root of $A'$ is a root of $A$. If the distinct roots of $A$ have
multiplicities $m_1,\ldots,m_q$, their multiplicities in $A'$ are
exactly $m_j-1$. They therefore account for $d-q$ roots of $A'$.
By hypothesis $A'$ has no zeros outside the zero set of $A$, so these
account for all $d-1$ zeros counted with multiplicity. Hence $q=1$.
\end{proof}

\begin{lemma}[The endpoints exclude pure powers]\label{lem:endpoints}
Let $f\in\Qb[x,y]$ have $A(t)=a(t-c)^d$, $a\ne0$, $d\ge1$.
Then $\HH$ is not rational.
\end{lemma}
\begin{proof}
Every nonzero value of $A$ is regular. At an actual endpoint $p$
with $A(p)=v\ne0$, the corresponding summand of the corrected leaf
germ extends holomorphically to $w=v$. Its Volterra integral vanishes
there, because the integrand is holomorphic and the integration path
shrinks to $p$. Its value is therefore, with the path orientation,
$E(p)/A'(p)$. Unramified preimages of a leaf have graph degree one,
so only actual endpoints can be active above this regular value.
By Lemma~\ref{lem:local-extraction}, rationality would force the sum
of these signed endpoint values to be zero.

If $A(0)\ne A(1)$, one endpoint has a nonzero value different from the
other endpoint value. Its nonzero leaf-germ value cannot cancel.
If the endpoint values agree, their common value $v$ cannot be zero,
since $A$ has just one distinct root. Rationality would give
\[
 \frac{\exp(A_1(0)/(Dv))}{A'(0)}
   =\frac{\exp(A_1(1)/(Dv))}{A'(1)}.
\]
The exponential of the algebraic difference of the two exponents is
therefore algebraic. Hermite--Lindemann forces that difference to be
zero; the equality then forces $A'(0)=A'(1)$. But
\[
 A'(0)=\frac{dv}{-c},\qquad A'(1)=\frac{dv}{1-c},
\]
which are unequal. When $d=1$, equal endpoint values are already
impossible. These cases exhaust the possibilities.
\end{proof}

\begin{proof}[Proof of Theorem~\ref{thm:nonflat}]
Suppose all positive factorial moments of a nonflat $f$ vanish.
Lemma~\ref{lem:specialization} gives such a polynomial with algebraic
coefficients. Its transform is $\HH=1/w$. If $d\ge2$,
Corollary~\ref{cor:annihilation} and Proposition~\ref{prop:arithmetic}
exclude every nonzero critical value. Lemma~\ref{lem:pure-power}
then makes $A$ a pure power. If $d=1$, it is already a translated
pure power. Lemma~\ref{lem:endpoints} contradicts rationality in
both cases.

Finally, if only finitely many moments of a nonflat $f$ were nonzero,
choose $N\ge1$ such that $\Lfac(f^n)=0$ for every $n\ge N$.
The nonzero polynomial $f^N$ has nonconstant projective top $A^N$
and every positive moment zero, contrary to the first assertion.
\end{proof}

\section{The flat case and the EIT input}\label{sec:flat}
The radial estimate below is independent of the monodromy argument.
It shows exactly where the EIT theorem enters.

\begin{lemma}[Uniform radial limit]\label{lem:radial-limit}
On every compact $t$-set on which $A$ does not vanish,
\begin{equation}\label{eq:radial-limit}
 \frac1{\Gamma(Dn+2)}\int_0^\infty e^{-r}rP(t,r)^n\dd r
 =A(t)^n\left[\exp\!\left(\frac{A_1(t)}{DA(t)}\right)+o(1)\right],
\end{equation}
with a uniform error inside the brackets.
\end{lemma}
\begin{proof}
For $\mathbf k=(k_1,\ldots,k_D)\in\mathbb N^D$, put
$K=\sum_jk_j$ and $J=\sum_jj k_j$. The multinomial expansion and
termwise radial integration give terms
\[
 A^{n-K}\prod_{j=1}^D\frac{A_j^{k_j}}{k_j!}\,C_{n,\mathbf k},
 \qquad
 C_{n,\mathbf k}=\frac{n!}{(n-K)!}
     \frac{\Gamma(Dn-J+2)}{\Gamma(Dn+2)}\quad(K\le n).
\]
Set $C_{n,\mathbf k}=0$ for $K>n$. The inequalities $K\le J\le DK$
show $0\le C_{n,\mathbf k}\le1$. Indeed, the smallest $K$ factors
in the denominator product dominate the numerator factors:
\[
 Dn-J+2+j\ \ge\ D(n-K)+2+j\ \ge\ n-K+1+j
 \qquad(0\le j<K).
\]
These denominator factors are present because $J\ge K$; all remaining
ones are at least one. The empty-product case $K=0$ is immediate.
For fixed $\mathbf k$,
\[
 C_{n,\mathbf k}\sim D^{-J}n^{K-J}.
\]
After division by $A^n$, only the terms $k_2=\cdots=k_D=0$ survive,
and their sum is $\exp(A_1/(DA))$.
On the compact set, all absolute terms are bounded by
\[
 \prod_{j=1}^D\frac{(b_j/a_*)^{k_j}}{k_j!},
 \qquad b_j=\max|A_j|,\quad a_* =\min|A|>0.
\]
This is summable over $\mathbb N^D$ and proves uniform dominated
convergence, including $D=1$.
\end{proof}

\begin{proof}[Proof of Corollary~\ref{cor:fc2}]
The nonflat case is Theorem~\ref{thm:nonflat}, and a nonzero constant
polynomial is immediate. Suppose $D\ge1$ and $A$ is constant.
It is nonzero, so scaling $f$ permits $A=1$. If all positive moments
vanish, specialize algebraically with this normalization retained.
Equation~\eqref{eq:radial-limit} gives
\[
 \lim_{n\to\infty}M_n=\int_0^1\exp\!\left(\frac{A_1(t)}D\right)\dd t.
\]
When $A_1$ is constant this limit is a nonzero exponential.
Otherwise Theorem~\ref{thm:eit} applies with $q=A_1$, $a=0$, $b=1$
and $\xi=1/D$. Its hypotheses hold over $\Qb$, and the amplitude is
exactly $\dd t$. The limit is again nonzero, contradicting $M_n=0$
for every $n\ge1$.
The infinitely-many assertion follows by applying the result to $f^N$
if a tail of its factorial moments were zero.
\end{proof}

\paragraph{An elementary flat subcase.}
The EIT input is unnecessary when $A$ is constant and $\deg A_1\le1$.
After normalizing $A=1$, retain the condition $\deg A_1\le1$ as rational
linear equations in the specialization argument. Write
$A_1(t)/D=\alpha t+\beta$ over $\Qb$. Then
\[
 \int_0^1e^{\alpha t+\beta}\dd t
 =\begin{cases}
 e^\beta,&\alpha=0,\\[2pt]
 e^\beta(e^\alpha-1)/\alpha,&\alpha\ne0.
 \end{cases}
\]
This is nonzero by Hermite--Lindemann~\cite{Baker}. Thus the same radial
limit proves nonvanishing in this subcase using published inputs only;
in particular it handles the flat case in total degree $D\le2$.
No rate of convergence beyond the uniform $o(1)$ in
Lemma~\ref{lem:radial-limit} is needed.

\paragraph{Research provenance and assistance.}
During the development of this argument, the \texttt{openai/math}
release suggested examining local jet separation, differential-system
multiplicity, and continuation of complete branch data
\cite{OAIJets,OAIMultiplicity,OAIContinuation}.
These were prompts for proof exploration, not theorems transferred into
the FC(2) argument. In particular, local \'etale rigidity does not imply
injectivity of a sum over sheets, and the cited billiard continuation
results are not applied here. The main proof uses none of the release's
results. Appendix~\ref{app:certificates} records the separate finite-jet
application of Moura's classical estimate, together with its formal
proof in the release. ChatGPT assisted with exploration, drafting,
symbolic checking and review; the author is responsible for the
mathematical claims and their attribution.

\appendix

\section{Exact continuation examples}\label{app:examples}
The examples in this appendix illustrate the continuation conventions.
None is a logical input to the main proof in
Sections~\ref{sec:kernel}--\ref{sec:flat}.

\subsection{The pole alone has the wrong jump}
For $f=3x+2y+1$, one has $D=1$ and $A(t)=2+t$. Its exponential
generating series is
\[
 \sum_{n\ge0}\frac{\Lfac(f^n)}{n!}z^n
       =\frac{e^z}{(1-3z)(1-2z)}.
\]
Integrating the series and substituting $z=1/w$ gives the exact identities
\[
 \HH(w)=\int_0^{1/w}\frac{e^z}{(1-3z)(1-2z)}\dd z,
 \qquad \HH'(w)=-\frac{e^{1/w}}{(w-3)(w-2)}.
\]
Counterclockwise continuation around $3$ therefore gives
$\Delta_3\HH=-2\pi i e^{1/3}$. In contrast, the pole integral alone is
\[
 C_0(w)=\int_2^3\frac{e^{1/u}}{w-u}\dd u,
 \qquad \Delta_3 C_0=-2\pi i e^{1/w}.
\]
The difference is a nonconstant analytic germ. For this family,
$\mathscr B(w,t)=e^{1/A(t)}/A(t)^2$, and the corrected branch germ
at a lower endpoint with value $v$ is
\[
 e^{1/w}+\int_v^w\frac{e^{1/u}}{u^2}\dd u=e^{1/v}.
\]
Thus the correction, while one local singular order softer, changes the
entire jump. A leading-pole identity cannot be substituted for
\eqref{eq:leaf-germ}.

\subsection{A coefficient that changes root sheets}
Take $D=2$ and $P=Ar^2+c$, where $c\ne0$ and $A=A(t)\ne0$.
The coefficient formula can be computed independently of radial
reversion. On writing $n=m+k$, the duplication formula gives
\[
 \binom{m+k}{k}\frac{\Gamma(2m+2)}{\Gamma(2m+2k+2)}
 =\frac{4^{-k}\Gamma(m+3/2)}{k!\,\Gamma(m+k+3/2)}.
\]
For $k\ge1$, use the beta integral for the ratio and sum the geometric
series over $m$. For large $|w|$ this gives
\begin{equation}\label{eq:quadratic-beta}
 \CC(w,t)=\frac1{w-A}
 +\int_0^1\frac{u^{1/2}}{w-Au}
       \sum_{k\ge1}\frac{(c/(4w))^k(1-u)^{k-1}}{k!(k-1)!}\dd u.
\end{equation}
All exchanges converge normally there. At the endpoint $u=1$, the
logarithmic coefficient is minus $1/A$ times the analytic numerator
evaluated at $u=w/A$. Reindexing $k=j+1$ yields
\begin{equation}\label{eq:quadratic-B}
 \mathscr B(w,t,\rho)
 =-\frac{c\rho}{4A^2}
   \sum_{j\ge0}\frac{[c(\rho^2-1)/(4A)]^j}{j!(j+1)!},
 \qquad \rho^2=A/w.
\end{equation}
This is an all-orders convergent identity. In particular,
$\mathscr B(A,t,1)=-c/(4A^2)$, agreeing with
\eqref{eq:global-B}.

The right side of \eqref{eq:quadratic-B} is odd in $\rho$. A fixed-$t$
loop around zero changes the sign of the continued distinguished
coefficient. When the loop returns to a neighborhood of the diagonal,
the canonical coefficient of the complete kernel is unchanged because
the singular root-filter summand has been reselected. This distinguishes
the two objects in Remark~\ref{rem:B-sheets} concretely.

\subsection{A moving endpoint with a changed lower root}
Let
\[
 P=(L(t)r+c)^D,\qquad L(t)=at+b,\qquad ac\ne0,\quad D\ge2.
\]
This is the radial polynomial of $[ax+b(x+y)+c]^D$.
Fix a lower endpoint $p$ with $L(p)\ne0$, and write
$s=L(t_i(w))$, so $s^D=w$. On the lifted integration path,
\[
 \Phi=r+c/L,\qquad \chi=Y-c/L,\qquad
 F=e^{c/L}(1-c/(LY)),\qquad \rho=L/s.
\]
Formula~\eqref{eq:global-B} now reads
\[
 \mathscr B\dd t=-\frac{\dd_t(e^{c/L(t)})}{Da\,s^{D-1}}.
\]
Since $A'(t_i)=Da\,s^{D-1}$, the full corrected expression is
\begin{equation}\label{eq:power-history}
 E(t_i(w))t_i'(w)+\int_p^{t_i(w)}\mathscr B\dd t
       =\frac{e^{c/L(p)}}{Da\,s^{D-1}}.
\end{equation}
For a positive loop around zero, $s$ changes to $\omega s$, where
$\omega=e^{2\pi i/D}$. The lower root parameter changes to
$\omega^{-1}\rho$, and \eqref{eq:power-history} changes by the factor
$\omega$. Resetting the lower root, or deleting the integral, would
both give the wrong continuation. The general transport proof retains
exactly these data.

\section{Finite certificates for branch densities}\label{app:certificates}
This appendix concerns auxiliary differential systems and is not an
input to the main proof. In particular, its branch-density tests do not
bound the first nonzero factorial moment. The analytic multiplicity
estimate is due to Moura~\cite{Moura,MouraHyper}. We prove the precise
polynomial-presentation, formal-power-series version below; the proof
is adapted from the version in~\cite{OAIMultiplicity}, which also gives
a Lean statement. The formal Wronskian step is explained in
Bostan--Dumas~\cite{BostanDumas}. These attributions concern the auxiliary
estimate, not a formal verification of FC(2).

\subsection{The classical multiplicity estimate}
\begin{theorem}[Moura's estimate in a polynomial presentation]
\label{thm:multiplicity}
Let $F$ have characteristic zero, $m,n\ge1$, $q\in F[z]$ with
$q(0)\ne0$ and $\deg q\le n$, and $P\in M_m(F[z])$ with
$\deg P_{ij}\le n-1$. Suppose
\[
 G\in\GL_m(F[[z]]),\qquad qG'=PG.
\]
For a constant row $u$ and constant column $v$, a nonzero series
$y=uGv$ satisfies
\begin{equation}\label{eq:moura-bound}
 \ord_0 y\le(n-1)\binom m2+m-1.
\end{equation}
\end{theorem}
\begin{proof}
Put $h=uG$, and let $V$ be the $F$-span of its component series.
Its dimension $k$ is between $1$ and $m$. Gaussian elimination on
successive coefficients gives a basis $b_1,\ldots,b_k$ of $V$ with
strictly increasing orders $a_1<\cdots<a_k$. The filtration by order
is eventually zero: in a finite-dimensional space any stabilizing
nonzero tail would lie in every power of $z$, which is impossible.
A nonzero element of $V$ has order at most $a_k$.

The Wronskian $W=\det(b_\ell^{(j)})_{0\le j<k,\,1\le\ell\le k}$
has order
\begin{equation}\label{eq:wronskian-order}
 \ord_0W=\sum_{i=1}^k a_i-\binom k2.
\end{equation}
Indeed, its leading coefficient after the powers of $z$ are factored
is the product of the leading coefficients of the $b_i$ times
$\det((a_i)_j)$, a nonzero Vandermonde determinant in characteristic
zero. Changing the basis of $V$ multiplies $W$ by a nonzero constant.

Define rows $r_0=u$, $r_{j+1}=r_j'+r_jP/q$. Then
$h^{(j)}=r_jG$. Writing $r_j=Z_j/q^j$ gives
\[
 Z_{j+1}=qZ_j'-j q'Z_j+Z_jP,\qquad \deg Z_j\le j(n-1).
\]
Let $R_*$ have rows $r_0,\ldots,r_{k-1}$ and put $S=\binom k2$.
Some $k$ components of $h$ form a basis of $V$, so $R_*G$ has a
nonzero $k$-minor by \eqref{eq:wronskian-order}. Thus $R_*$ also has
such a minor. Each $k$-minor of $R_*$ has denominator $q^S$ and a
polynomial numerator of degree at most $(n-1)S$. Since $q$ is a unit
at zero, the minimum order $\delta$ of its nonzero $k$-minors obeys
$0\le\delta\le(n-1)S$.

Multiplication by $G$ preserves the ideal of $k$-minors in $F[[z]]$:
apply Cauchy--Binet first with $G$ and then with $G^{-1}$. Hence the
minimum order for the minors of $R_*G$ is also $\delta$. A nonzero
minor of that order is a Wronskian of $k$ components of $h$; those
components form a basis of $V$. Equation~\eqref{eq:wronskian-order}
therefore gives
\[
 \delta=\sum_{i=1}^k(a_i-(i-1)),\qquad
 \ord_0y\le a_k\le\delta+k-1
     \le(n-1)\binom k2+k-1.
\]
Each summand in the first expression is nonnegative. Finally $k\le m$
proves \eqref{eq:moura-bound}.
\end{proof}
\paragraph{Comparison with the analytic statement.}
Over $\C$, Moura's formulation~\cite[Section~2, Theorem~1]{MouraHyper}
uses the total pole order $\sigma$ of the matrix differential form
$\Omega=(P/q)\dd z$ on the projective line. If $P\ne0$, put
$s=\deg q$ and $p=\max_{P_{ij}\ne0}\deg P_{ij}$. Finite poles contribute
at most $s$, and the pole order at infinity is at most
$\max(0,p-s+2)$. Therefore
\[
 \sigma\le s+\max(0,p-s+2)=\max(s,p+2)\le n+1.
\]
Moura's expression $m-1+(\sigma-2)\binom m2$ gives
\eqref{eq:moura-bound}. If $P=0$, the solution matrix is constant and
the bound is immediate. This comparison does not assume $\deg q=n$;
the direct proof above applies over every characteristic-zero field.
The corresponding pinned Lean declaration is
\texttt{OAI.NCHitting.}\allowbreak\texttt{system\_multiplicity}.
Citing or reading that declaration does not claim that the present
FC(2) paper, or the density application below, has been formally verified.

\subsection{An explicit rational connection}
Fix $A,B\in\C[t]$, $\deg A=d\ge1$, and an integer $D\ge1$.
On a disk avoiding zero and the critical values of $A$, put
\[
 h_i(w)=\frac{\exp(B(t_i(w))/(Dw))}{A'(t_i(w))},\qquad
 \Psi_{ji}(w)=t_i(w)^j h_i(w)\quad(0\le j<d).
\]
The determinant of $\Psi$ is a Vandermonde determinant times nonzero
column factors, so it is invertible. In $\C(w)[t]/(A(t)-w)$, reduce
\begin{equation}\label{eq:density-connection}
 \frac{j t^{j-1}}{A'}-\frac{t^jA''}{(A')^2}
 +\frac{t^jB'}{DwA'}-\frac{t^jB}{Dw^2}
       =\sum_{k=0}^{d-1}R_{jk}(w)t^k.
\end{equation}
The first term is understood as zero for $j=0$.
Differentiating $t_i'=1/A'(t_i)$ and $h_i$ proves the exact system
$\Psi'=R\Psi$.

Let $\Delta(w)=\Res_t(A(t)-w,A'(t))$. It is a nonzero polynomial.
The resultant identity provides a representative of $1/A'$ whose
only parameter denominator is $\Delta$. Reduction modulo $A-w$
introduces only constant leading-coefficient denominators. Consequently
\[
 q(w)=Dw^2\Delta(w)^2,\qquad P(w)=q(w)R(w)
\]
has polynomial entries. Define
\[
 n=\max\{1,\deg q,1+\max_{P_{ij}\ne0}\deg P_{ij}\},\qquad
 J=(n-1)\binom d2+d-1,
\]
omitting the last entry of the maximum when $P=0$.

\begin{corollary}[Finite test for an identically zero density]
\label{cor:finite-jets}
For a constant coefficient vector $c$, a point $w_0$ with $q(w_0)\ne0$
(equivalently $w_0\notin\{0\}\cup\Crit(A)$ for the displayed $q$),
and $h(w)=\sum_i c_i h_i(w)$ on a sufficiently small disk about $w_0$,
one has
\[
 h\equiv0\quad\Longleftrightarrow\quad
 h^{(j)}(w_0)=0\quad(0\le j\le J).
\]
\end{corollary}
\begin{proof}
Translate $w_0$ to zero and normalize the fundamental matrix to
$G(w)=\Psi(w)\Psi(w_0)^{-1}$. Then
$h=e_0^{\mathsf T}G(\Psi(w_0)c)$ is a constant scalar output.
The denominator $q$ is nonzero at $w_0$. Apply
Theorem~\ref{thm:multiplicity} contrapositively.
\end{proof}
For algebraic $A,B,w_0,c$, the derivatives are finite algebraic linear
combinations of exponentials of algebraic numbers. They can be tested
exactly by grouping equal exponents and using Lindemann--Weierstrass.
This certifies a particular density identity. It neither makes every
sheet sum nonzero nor bounds the first nonzero moment of an arbitrary
factorial polynomial.

For example, $A=t^2+1$, $B=t$ and $D=2$ give
\[
 R=\begin{pmatrix}
 -\dfrac1{2(w-1)}&\dfrac{2-w}{4w^2(w-1)}\\[5pt]
 \dfrac{2-w}{4w^2}&0
 \end{pmatrix}.
\]
One may take $q=4w^2(w-1)$, $n=3$ and $J=3$.
The first two observation rows $e_0^{\mathsf T}$ and
$e_0^{\mathsf T}R$ are independent over $\C(w)$.
By contrast, for $A=t(1-t)$ and constant $B=c$,
\[
 \frac{e^{c/(2w)}}{A'(t_+(w))}
 +\frac{e^{c/(2w)}}{A'(t_-(w))}=0,
 \qquad t_\pm(w)=\frac{1\pm\sqrt{1-4w}}2.
\]
Both summands are nonzero. All jets of their sum vanish, even on the
\'etale locus. The actual interval uses opposite orientations on these
two branches; this is why the signed path data in the main proof
cannot be replaced by local jet rigidity.

\subsection{A different, conditional certificate for a complete system}
If an exact full transform is embedded in a rational system $Y'=MY$
and $h=\ell_0Y$, define observation rows
$\ell_{j+1}=\ell_j'+\ell_jM$. Then $h^{(j)}=\ell_jY$.
If a coordinate row $e_r^{\mathsf T}$ lies in the rational span of
these rows and $Y_r\not\equiv0$, one necessarily has $h\not\equiv0$.
Indeed, $h\equiv0$ would force $Y_r\equiv0$ by the displayed row relation.
This criterion requires actual observation of a nonzero coordinate;
finiteness of the system alone is not sufficient.

For a nonconstant linear polynomial $f=ax+by+c$, let
$q=(w-a)(w-b)$ and take
\[
 Y=(\HH,e^{c/w},1)^{\mathsf T},\qquad
 M=\begin{pmatrix}0&-1/q&0\\0&-c/w^2&0\\0&0&0\end{pmatrix},
 \qquad \ell_0=(1,0,-1/w).
\]
The system follows by integrating the exponential generating series,
and $h=\HH-1/w$ is the exact positive-moment output. Direct differentiation
gives
\[
 \det\begin{pmatrix}\ell_0\\\ell_1\\\ell_2\end{pmatrix}
 =\frac{-(a+b+c)w^2+(2ab+c(a+b))w-cab}{q^2w^4}.
\]
Its numerator is identically zero only if $a=b=c=0$:
from $c=-a-b$, the middle coefficient gives $a^2+b^2=0$,
and the last gives $ab(a+b)=0$, forcing $a=b=0$.
Thus in the nonconstant case all coordinates, including $1$, are
observed and $h$ cannot vanish identically. No analogous universal
observation assertion in higher degree is needed or claimed here.

\section{A collision with incomparable decompositions}\label{app:collision}
Here is a concrete instance of Proposition~\ref{prop:arithmetic}.
Let $D\ge6$ and suppose
\[
 A(t)=T_6(2t-1)-1,\qquad
 T_6(z)=32z^6-48z^4+18z^2-1.
\]
It admits both $T_6=T_2\circ T_3$ and $T_6=T_3\circ T_2$, where
$T_2(z)=2z^2-1$ and $T_3(z)=4z^3-3z$. The corresponding block sizes
are incomparable, and $A(0)=A(1)=0$.
The three points over the nonzero critical value $-2$ are
\[
 c_-=\frac{2-\sqrt3}{4},\qquad c_0=\frac12,\qquad
 c_+=\frac{2+\sqrt3}{4}.
\]
They have local degree two, with quadratic coefficients
$A''(c)/2$ equal to $288,72,288$, respectively. The positive
critical-point weights in \eqref{eq:positive-pairing} are therefore
$1/576,1/144,1/576$. Rationality of $\HH$ would force
\[
 e^{-A_1(c_-)/(2D)}+4e^{-A_1(1/2)/(2D)}
                  +e^{-A_1(c_+)/(2D)}=0.
\]
For algebraic layers this contradicts Lindemann--Weierstrass after equal
exponents are grouped. There are no restrictions on the subsequent
layers. Algebraic specialization gives the corresponding nonvanishing
statement with arbitrary complex coefficients. The computation illustrates
why neither colliding critical points, incomparable block systems, nor
zero endpoint values are omitted by the proof.

\paragraph{Reproducibility.}
The source bundle includes an exact $\Q(\sqrt3)$ calculation of the
identities in this appendix, and symbolic checks of the rational
connection, observation determinant, root filter, radial normalization,
and branch-change examples. These finite checks are separate from the
analytic and representation-theoretic proofs.
\small
\begin{thebibliography}{99}
\bibitem{AmazingImage}
A.~van~den~Essen, \emph{The Amazing Image Conjecture},
arXiv:1006.5801v1 (2010).
\href{https://arxiv.org/abs/1006.5801v1}{arXiv:1006.5801v1}.

\bibitem{Baker}
A.~Baker, \emph{Transcendental Number Theory}, Cambridge Mathematical
Library, Cambridge University Press, 1990, Chapter~1.

\bibitem{BostanDumas}
A.~Bostan and P.~Dumas, \emph{Wronskians and linear independence},
Amer. Math. Monthly \textbf{117} (2010), no.~8, 722--727.
\href{https://arxiv.org/abs/1301.6598}{arXiv:1301.6598}.

\bibitem{FC}
A.~van~den~Essen, D.~Wright and W.~Zhao,
\emph{On the Image Conjecture}, J. Algebra \textbf{340} (2011),
211--224.
\href{https://doi.org/10.1016/j.jalgebra.2011.04.036}{doi:10.1016/j.jalgebra.2011.04.036}.
\href{https://arxiv.org/abs/1008.3962v2}{arXiv:1008.3962v2}.

\bibitem{LiuSun}
D.~Liu and X.~Sun, \emph{The factorial conjecture and images of locally
nilpotent derivations}, Bull. Aust. Math. Soc. \textbf{101} (2020),
71--79.
\href{https://doi.org/10.1017/S0004972719000546}{doi:10.1017/S0004972719000546}.

\bibitem{PM}
F.~Pakovich and M.~Muzychuk, \emph{Solution of the polynomial moment
problem}, Proc. Lond. Math. Soc. (3) \textbf{99} (2009), 633--657.
\href{https://doi.org/10.1112/plms/pdp010}{doi:10.1112/plms/pdp010}.
\href{https://arxiv.org/abs/0710.4085v2}{arXiv:0710.4085v2}.
Numbering refers to arXiv~v2: Sections~2.1--2.2 and~4.1,
Theorem~3.1, and Propositions~2.4, 2.6 and~4.1.

\bibitem{EIT}
C.~D.~Long and A.-A.~Le~Blanc, \emph{Algebraic Exponential Integrals:
Transcendence and Linear Relations}, EIT~I, research manuscript, 2026,
Corollary~1.4 (source label \texttt{cor:single}).
\href{https://github.com/long-mathematics/exponential-integral-theorem/blob/22d98c4f27e606b5b25fa673c33f9c56c53cc86b/papers/01-exponential-integral-theorem/exponential-integral-theorem.tex}{Pinned LaTeX source},
commit \texttt{22d98c4f27e606b5b25fa673c33f9c56c53cc86b},
consulted 7 October 2026; the separate input to Theorem~\ref{thm:eit}.

\bibitem{Moura}
C.~Moura, \emph{Bounds for the vanishing order for solutions of a linear
differential system}, J. Dynam. Control Systems \textbf{9} (2003),
no.~1, 73--88.
\href{https://doi.org/10.1023/A:1022155217548}{doi:10.1023/A:1022155217548}.

\bibitem{MouraHyper}
C.~Moura, \emph{On the multiplicity of the hyperelliptic integrals},
preprint, 2003, Section~2, Theorem~1.
\href{https://arxiv.org/abs/math/0312323}{arXiv:math/0312323}.
Analytic total-pole-order formulation; see Appendix~\ref{app:certificates}.

\bibitem{OAIMultiplicity}
OpenAI, \emph{One Rational Matrix Hitting Point for Noncommutative
Formulas}, manuscript dated 24 September 2026, section
\emph{A bound on the order of vanishing}, in \texttt{openai/math}.
\href{https://github.com/openai/math/blob/adc7f1241b42e322a6451854ab7e4b4c146bf78a/preprints/One-Rational-Matrix-Hitting-Point-for-Noncommutative-Formulas-September-24-2026/build/multiplicity.tex}{Manuscript proof};
\href{https://github.com/openai/math/blob/adc7f1241b42e322a6451854ab7e4b4c146bf78a/lean/OAI/Algebra/FormulaHitting/Multiplicity.lean}{Lean source},
\texttt{OAI.NCHitting.}\allowbreak\texttt{system\_multiplicity}.
The manuscript attributes the estimate to Moura and its formal
Wronskian ingredient to Bostan--Dumas. Its prose uses $\deg q=n$;
the cited Lean statement allows $\deg q\le n$. Sources pinned to
commit \texttt{adc7f1241b42e322a6451854ab7e4b4c146bf78a},
consulted 7 October 2026.

\bibitem{OAIJets}
OpenAI, \emph{\'Etale analytic jet lemmas},
\href{https://github.com/openai/math/blob/adc7f1241b42e322a6451854ab7e4b4c146bf78a/lean/OAI/AlgebraicGeometry/Seshadri/Jets/EtaleAnalytic.lean}{\texttt{EtaleAnalytic.lean}},
\texttt{etale\_tangent\_separates}, in \texttt{openai/math},
commit \texttt{adc7f1241b42e322a6451854ab7e4b4c146bf78a},
consulted 7 October 2026. Cited for research provenance, not as an
input to the main proof.

\bibitem{OAIContinuation}
OpenAI, \emph{Rigidity of Smooth Billiards with a Continuous Caustic
Collar}, manuscript dated 24 September 2026, section
\emph{Endpoint determination and meromorphic continuation}.
\href{https://github.com/openai/math/blob/adc7f1241b42e322a6451854ab7e4b4c146bf78a/preprints/Rigidity-of-Smooth-Billiards-with-a-Continuous-Caustic-Collar-September-24-2026/build/sections/04-continuation.tex}{Pinned section source},
\texttt{openai/math}, commit
\texttt{adc7f1241b42e322a6451854ab7e4b4c146bf78a},
consulted 7 October 2026. Cited for research provenance only; the
billiard hypotheses are not asserted for the present integrals.
\end{thebibliography}
\end{document}
